\documentclass[11pt]{article}
\usepackage[margin=1in]{geometry}
\usepackage{amsmath,amssymb,amsthm}
\usepackage{enumitem}
\usepackage{mathtools}
\usepackage{hyperref}
\usepackage{pgffor}

\newtheorem*{theorem*}{Theorem}

\title{MIT 6.1200J Problem Set 2 – Solutions}
\author{Nouman Rahman}
\date{1st September 2026}

\newcommand{\N}{\mathbb{N}}
\newcommand{\R}{\mathbb{R}}
\newcommand{\Z}{\mathbb{Z}}
\newcommand{\TODO}[1]{\emph{[TODO: #1]}}
\newcommand{\partlabel}[1]{%
    \noindent\makebox[2em][l]{\textbf{(#1)}}%
}

\begin{document}
\maketitle

\section*{Problem 1: Outlining Proofs}

\partlabel{a}
Theorem: ``Every natural number is tall or wide, but not both.'' (Proof method: proof by cases.)

Writing $T(n)$ for ``$n$ is tall'' and $W(n)$ for ``$n$ is wide,'' the theorem says
\[
\forall n \in \N.\ \big(T(n) \oplus W(n)\big),
\]
where $\oplus$ is exclusive or. We outline a proof by cases on whether $n$ is tall or wide (this case split itself needs $T(n) \vee W(n)$ to be already known, e.g.\ from an earlier result — we flag this as a TODO as well).

\begin{quote}
\textit{Proof Outline.} Let $n \in \N$ be arbitrary. \TODO{show $T(n) \vee W(n)$ (needed so the two cases below are exhaustive)}. We consider two cases.

\textbf{Case 1: $T(n)$ holds.} We must show $T(n) \oplus W(n)$, i.e.\ (since $T(n)$ is true) that $W(n)$ is false. \TODO{show $T(n) \rightarrow \neg W(n)$}.

\textbf{Case 2: $W(n)$ holds.} Symmetrically, we must show $T(n)$ is false. \TODO{show $W(n) \rightarrow \neg T(n)$}.

In both cases $T(n) \oplus W(n)$ holds, completing the proof by cases.
\end{quote}

\bigskip

\partlabel{b}
Theorem: ``There exists a happy natural number between $10$ and $1000$ that has no happy proper divisors.'' (Proof method: choose $25$ as the example.)

Writing $H(n)$ for ``$n$ is happy'' and $D(n)$ for the set of proper divisors of $n$, the theorem says
\[
\exists n \in \N.\ \Big( 10 \le n \le 1000 \ \wedge\ H(n) \ \wedge\ \forall d.\big(d \mid n \wedge d < n \rightarrow \neg H(d)\big) \Big).
\]

\begin{quote}
\textit{Proof Outline.} Take $n = 25$.

$10 \le 25 \le 1000$: true by inspection.

\TODO{show $H(25)$, i.e.\ that $25$ is happy}.

The proper divisors of $25$ are $1$ and $5$. \TODO{show $\neg H(1)$ and $\neg H(5)$}.

Since $25$ is in range, happy, and has no happy proper divisor, $n = 25$ witnesses the existential claim.
\end{quote}

\bigskip
\partlabel{c}
Theorem: ``There are no devious natural numbers.'' (Proof method: regular induction, with base cases $0$ and $1$.)

Writing $D(n)$ for ``$n$ is devious,'' the theorem says $\forall n \in \N.\ \neg D(n)$. Let
\[
P(n) := \neg D(n).
\]
Two base cases are called for, which suggests the (unstated) recursive structure of ``devious'' relates $n$ to its two predecessors $n-1, n-2$ (as in Problem 2). We outline accordingly:

\begin{quote}
\textit{Proof Outline.} We prove $P(n)$ for all $n \in \N$ by induction on $n$, with two base cases and a two-term inductive step.

\textbf{Base case $n = 0$:} \TODO{show $P(0)$, i.e.\ $\neg D(0)$}.

\textbf{Base case $n = 1$:} \TODO{show $P(1)$, i.e.\ $\neg D(1)$}.

\textbf{Inductive step:} Let $k \ge 1$, and assume $P(k-1)$ and $P(k)$ (i.e.\ $\neg D(k-1)$ and $\neg D(k)$) both hold. \TODO{show $P(k+1)$, i.e.\ $\neg D(k+1)$, using the definition of devious and the two inductive hypotheses}.

By induction (using both base cases as anchors), $P(n)$ holds for all $n \in \N$.
\end{quote}

\newpage
\section*{Problem 2: Spotting Errors}

Recall $x_1 = 4$, $x_2 = 13$, and $x_i = 6x_{i-1} - 8x_{i-2}$ for $i \ge 3$, so $(x_1,x_2,x_3,x_4,\dots) = (4,13,46,172,\dots)$. Goal: $x_i > 0$ for all $i \ge 1$.

\medskip
\partlabel{a}
Befuddled Brynmor's mistake.

Brynmor uses strong induction on $Q(i) := $ ``$x_i \ge 2^i$'', and in the inductive step writes
\[
x_{i+1} = 6x_i - 8x_{i-1} \ \ge\ 6\cdot 2^i - 8 \cdot 2^{i-1},
\]
substituting the lower bounds $x_i \ge 2^i$ and $x_{i-1} \ge 2^{i-1}$ directly into the expression.

\emph{This step is invalid.} The term $6x_i$ has a positive coefficient, so substituting the lower bound $x_i \ge 2^i$ correctly gives $6x_i \ge 6\cdot 2^i$. But the term $-8x_{i-1}$ has a \emph{negative} coefficient: to lower-bound $-8x_{i-1}$ we need an \emph{upper} bound on $x_{i-1}$, not a lower bound. Knowing $x_{i-1} \ge 2^{i-1}$ only tells us $-8x_{i-1} \le -8\cdot 2^{i-1}$ — the inequality goes the wrong way to conclude anything about $x_{i+1}$ being large. So the displayed inequality $x_{i+1} \ge 6\cdot 2^i - 8\cdot 2^{i-1}$ does not actually follow from the inductive hypothesis, and the rest of the argument is unjustified.

\medskip
\partlabel{b}
Erroneous Erik's mistake.

Erik uses (regular) induction on $R(i) := $ ``$x_i \ge 3x_{i-1}$'' for $i \ge 2$. In the inductive step, after correctly deriving $x_{i+1} \ge \tfrac{10}{3}x_i$ from the inductive hypothesis $R(i)$, he writes: ``Since $x_i > 0$, $\tfrac{10}{3}x_i > 3x_i$'' in order to conclude $x_{i+1} \ge 3x_i$, i.e.\ $R(i+1)$.

\emph{This step is circular.} The inductive step for $R(i+1)$ is only allowed to use the inductive hypothesis $R(i)$ (namely $x_i \ge 3x_{i-1}$) — it does \emph{not} give us $x_i > 0$ for free. Positivity of $x_i$ is exactly the fact the whole exercise is trying to establish, and it is only derived \emph{afterward}, by combining $R(i)$ with $x_1 > 0$. Using ``$x_i > 0$'' inside the proof of $R(i)\Rightarrow R(i+1)$, before positivity has been established, assumes the very conclusion the argument is building toward.

\medskip
\partlabel{c}
Overzealous Zach's mistake.

Zach uses strong induction on $S(i) := $ ``$x_i = 4^i$'', with a single base case $S(1)$, and an inductive step that (for general $i\ge 1$) invokes both $S(i-1)$ and $S(i)$ to derive $S(i+1)$ via the recurrence $x_{i+1} = 6x_i - 8x_{i-1}$.

\emph{This step fails already at $i = 1$.} To derive $S(2)$ this way, the inductive step needs $S(0)$, i.e.\ a value $x_0 = 4^0 = 1$ — but $x_0$ is not defined at all ($x_1, x_2$ are the given base values, and the recurrence $x_i = 6x_{i-1}-8x_{i-2}$ is only stipulated for $i \ge 3$). So the strong induction step's use of the recurrence to compute $x_2$ from $x_1, x_0$ is illegitimate on two counts: $x_0$ doesn't exist, and even $x_2$ itself is a \emph{given} base value, not something produced by the recurrence. Consistent with this, the claim $S(2)$ is simply false: $x_2 = 13 \ne 16 = 4^2$. So the strong induction breaks down exactly at the step from $i=1$ to $i=2$.

\bigskip
\partlabel{d}
Correct proof that $x_i > 0$ for all $i \ge 1$.

Following the hint, we repair Erik's argument by folding positivity directly into the inductive hypothesis, so it no longer needs to be smuggled in mid-proof. Let
\[
P(i) := \text{``}x_i > 2x_{i-1}\text{''}, \qquad \text{for } i \ge 2.
\]

\begin{proof}
We prove $P(i)$ holds for all $i \ge 2$ by induction on $i$.

\textbf{Base case ($i=2$):} $x_2 = 13$ and $2x_1 = 8$, and indeed $13 > 8$, so $P(2)$ holds.

\textbf{Inductive step:} Let $i \ge 2$ and assume $P(i)$, i.e.\ $x_i > 2x_{i-1}$. We show $P(i+1)$, i.e.\ $x_{i+1} > 2x_i$.

From $x_i > 2x_{i-1}$ we get $x_{i-1} < \tfrac{1}{2}x_i$. Multiplying both sides by $-8$ (which flips the inequality, since $-8 < 0$):
\[
-8x_{i-1} > -8\cdot\tfrac{1}{2}x_i = -4x_i.
\]

Now use the recurrence $x_{i+1} = 6x_i - 8x_{i-1}$:
\[
x_{i+1} = 6x_i - 8x_{i-1} \ >\ 6x_i - 4x_i = 2x_i.
\]

So $x_{i+1} > 2x_i$, which is $P(i+1)$.

By induction, $P(i)$ holds for all $i \ge 2$, i.e.\
\[
x_i > 2x_{i-1} \qquad \text{for all } i \ge 2. \tag{$\ast$}
\]

Finally, we conclude $x_i > 0$ for all $i \ge 1$ by a second, simple induction, using $(\ast)$:
\begin{itemize}
\item $x_1 = 4 > 0$.
\item If $x_{i-1} > 0$ for some $i \ge 2$, then by $(\ast)$, $x_i > 2x_{i-1} > 2\cdot 0 = 0$.
\end{itemize}
Hence $x_i > 0$ for every $i \ge 1$, as required.
\end{proof}

\newpage
\section*{Problem 3: Integer Multiplication from Bit Shifts}

State machine on triples $(r,s,a) \in \N^3$, start state $(x,y,0)$, transitions (defined only when $s > 0$):
\[
(r,s,a) \mapsto
\begin{cases}
(2r,\ s/2,\ a) & \text{if } s \text{ is even (and } s>0\text{)} \\
(2r,\ (s-1)/2,\ a+r) & \text{if } s \text{ is odd}
\end{cases}
\]

\medskip
\partlabel{a}
Tracing from $(5,22,0)$ (target: $5\cdot 22 = 110$):
\[
(5,22,0) \to (10,11,0) \to (20,5,10) \to (40,2,30) \to (80,1,30) \to (160,0,110).
\]
No transition applies once $s=0$, so this is the final state, with $a_f = 110 = 5 \cdot 22$. It correctly computes the product.

\medskip
\partlabel{b}
\textbf{Invariant.} Let $P(r,s,a) := $ ``$rs + a = xy$.'' We show $P$ is an invariant of the state machine: true at the start state, and preserved by every transition.

\begin{proof}
\textbf{Holds initially:} $P(x,y,0)$ says $xy + 0 = xy$, which is true.

\textbf{Preserved by transitions:} Suppose $P(r,s,a)$ holds, i.e.\ $rs+a = xy$, and $s > 0$ so a transition applies.

\emph{Case $s$ even:} $(r,s,a) \mapsto (2r,\ s/2,\ a)$. Then
\[
(2r)\left(\frac{s}{2}\right) + a = rs + a = xy,
\]
so $P(2r,s/2,a)$ holds.

\emph{Case $s$ odd:} $(r,s,a) \mapsto (2r,\ (s-1)/2,\ a+r)$. Then
\[
(2r)\left(\frac{s-1}{2}\right) + (a+r) = r(s-1) + a + r = rs - r + a + r = rs+a = xy,
\]
so $P(2r,(s-1)/2,a+r)$ holds.

In both cases, the successor state also satisfies $P$. Since $P$ holds at the start state and is preserved by every transition, $P$ is an invariant.
\end{proof}

\medskip
\partlabel{c}
\textbf{Final states.} Transitions are defined exactly when $s > 0$, so a state has no outgoing transition (is final) precisely when $s = 0$. Thus the final states are exactly the triples of the form $(r, 0, a)$ for $r, a \in \N$.

\medskip
\partlabel{d}
\textbf{Partial correctness.} Let $(r_f, s_f, a_f)$ be any reachable final state. By part (c), $s_f = 0$. Since $P$ is an invariant (part (b)) and holds at the start state, it holds at every reachable state, including this one:
\[
P(r_f, 0, a_f):\qquad r_f \cdot 0 + a_f = xy \quad\Longrightarrow\quad a_f = xy.
\]
So in any reachable final state, $a_f = xy$, as required.

\medskip
\partlabel{e}
\textbf{Termination.} Let $\mathrm{bitcount}(n) := \lceil \log_2(n+1) \rceil$ for $n \in \N$. We show $\mathrm{bitcount}(s)$ strictly decreases — in fact, decreases by exactly $1$ — with every transition.

It is a standard fact that for $n \ge 1$, $\mathrm{bitcount}(n)$ is the unique integer $k$ such that $2^{k-1} \le n < 2^k$ (i.e.\ the number of bits in $n$'s binary representation), and $\mathrm{bitcount}(0) = 0$.

\begin{proof}
Let $(r,s,a)$ be any state with $s > 0$, and suppose $\mathrm{bitcount}(s) = k$, so $2^{k-1} \le s < 2^k$.

\textbf{Case $s$ odd:} the transition sends $s \mapsto \dfrac{s-1}{2}$. Write $m = \dfrac{s-1}{2}$, so $s = 2m+1$.
From $2^{k-1} \le s < 2^k$ we get $2^{k-1} \le 2m+1 < 2^k$, i.e.\ $2^{k-1}-1 \le 2m < 2^k - 1$; since $2m$ is even and $2^{k-1}-1$ is odd (for $k \ge 1$), this sharpens to $2^{k-1} \le 2m \le 2^k - 2$, i.e.\ $2^{k-2} \le m \le 2^{k-1}-1 < 2^{k-1}$. If $k \ge 2$ this says $2^{(k-1)-1} \le m < 2^{k-1}$, so $\mathrm{bitcount}(m) = k-1$. (If $k=1$, then $s=1$, $m=0$, and $\mathrm{bitcount}(0)=0=k-1$ by definition.) Either way,
\[
\mathrm{bitcount}\!\left(\frac{s-1}{2}\right) = k - 1 = \mathrm{bitcount}(s) - 1.
\]

\textbf{Case $s$ even:} the transition sends $s \mapsto \dfrac{s}{2}$. Write $m = \dfrac{s}{2}$, so $s = 2m$, and (since $s>0$) $m \ge 1$.
From $2^{k-1} \le s < 2^k$ we get $2^{k-1} \le 2m < 2^k$, i.e.\ $2^{k-2} \le m < 2^{k-1}$, so $\mathrm{bitcount}(m) = k-1$. Hence
\[
\mathrm{bitcount}\!\left(\frac{s}{2}\right) = k-1 = \mathrm{bitcount}(s) - 1.
\]

In both cases, the new value of $s$ has $\mathrm{bitcount}$ exactly one less than before. In particular, $\mathrm{bitcount}(s)$ is a strictly decreasing, non-negative-integer-valued derived variable along any run of the machine. Since it cannot decrease forever (it is a non-negative integer), the machine cannot run forever: it must reach a state with $s = 0$ (no further transition possible), i.e.\ the algorithm always terminates.
\end{proof}

\medskip
\partlabel{f}
\textbf{Time bound.} By part (e), $\mathrm{bitcount}(s)$ decreases by exactly $1$ at every step, starting from $\mathrm{bitcount}(y)$ at the initial state $(x,y,0)$ and ending at $\mathrm{bitcount}(0) = 0$ at the final state. Hence, the number of transitions taken is exactly
\[
\mathrm{bitcount}(y) = \lceil \log_2(y+1) \rceil,
\]
so from any starting state $(x,y,0)$, the algorithm terminates after at most $\lceil \log_2(y+1)\rceil$ steps.

\newpage
\section*{Problem 4: Invariant Enlightenment}

\end{document}
